\begin{tikzpicture}[font=\sffamily\small]
% ---------------- the Linux side
\node[user, text width=32mm] (pub) at (0,1.4) {mosquitto\_pub\\{\scriptsize lab/p16/led/red}};
\node[user, text width=32mm] (broker) at (0,0) {mosquitto broker\\{\scriptsize port 1883, from P12}};
\node[kern, text width=32mm] (net) at (0,-1.6) {eth0 / usb0\\{\scriptsize Ethernet and the LTE path}};
\node[hw, text width=32mm] (hdr) at (0,-3.2) {Pi 4 40-pin header\\{\scriptsize empty in this lab}};
% ---------------- the link
\node[ext, text width=28mm] (lan) at (4.6,0) {LAN or the Pi\\access point};
% ---------------- the ESP32 side
\node[fw, text width=30mm] (esp) at (9.2,0) {ESP32 firmware\\{\scriptsize Wi-Fi STA, MQTT client, three outputs}};
\node[hw, text width=26mm] (leds) at (13.4,0) {R/Y/G LEDs\\{\scriptsize low = lit}};
% ---------------- domains
\begin{scope}[on background layer]
\node[layer, fit=(pub)(broker)(net)(hdr), inner ysep=9pt, inner xsep=5pt,
  label={[layerlabel, anchor=south west]north west:Linux keeps storage, time, Ethernet and LTE}] {};
\node[layer, fit=(esp)(leds), inner ysep=9pt, inner xsep=5pt,
  label={[layerlabel, anchor=south west]north west:the ESP32 keeps the 2.4 GHz radio and three pins}] {};
\end{scope}
% ---------------- links
\draw[flow] (pub) -- (broker);
\draw[bus] (broker) -- node[lbl, above]{TCP 1883} (lan);
\draw[bus] (lan) -- node[lbl, above]{2.4 GHz} (esp);
\draw[flow] (esp) -- node[lbl, above]{pin} (leds);
\draw[bus] (net.west) -- ++(-6mm,0) |- node[lbl, left, pos=0.25]{WAN} (broker.west);
% ---------------- the wire that is deliberately absent
\draw[red!70!black, dashed, line width=0.8pt] (hdr.south) -- (0,-4.6) -- (13.4,-4.6) -- (leds.south);
\draw[red!70!black, line width=1pt] (6.5,-4.85) -- (6.9,-4.35);
\draw[red!70!black, line width=1pt] (6.5,-4.35) -- (6.9,-4.85);
\node[font=\sffamily\scriptsize, text=red!70!black, anchor=north] at (6.7,-4.95)
  {the Pi does not bit-bang these LEDs};
\end{tikzpicture}
